% ============================================================
% RESUMEN
% ============================================================

\begin{abstract}

En este trabajo se realiza un análisis espectral de señales de audio mediante la Transformada Rápida de Fourier (FFT), utilizando la herramienta desarrollada en el proyecto \texttt{WORK\_THREE} (paquetes \texttt{signal\_lab} y \texttt{animal\_analyzer}). Se estudiaron tres categorías de señales: sonidos producidos por animales, sonidos correspondientes a instrumentos musicales y voces humanas. Las grabaciones de animales (gato, perro, gallo) e instrumentos (contrabajo, piano, flauta) corresponden a audio real de dominio público/Creative Commons descargado de Wikimedia Commons; las dos señales de voz se sintetizaron localmente con Windows SAPI ante la falta de grabaciones humanas, y se documentan como tales. Cada grabación se cargó a su frecuencia de muestreo nativa y se acotó a los primeros tres segundos, siguiendo la ventana de análisis fija (\texttt{ANALYSIS\_WINDOW\_S = 3.0~s}) implementada en \texttt{animal\_analyzer/core.py}. Para cada señal se obtuvo su representación en el dominio temporal y frecuencial, además de calcular la media, la desviación estándar y la frecuencia dominante. Adicionalmente, se introdujo una métrica cuantitativa de separabilidad estadística entre pares de señales, definida como la razón entre la diferencia de frecuencias dominantes y la suma de sus desviaciones estándar. Los siete pares evaluados resultaron, sin excepción, claramente separables según esta métrica; sin embargo, esa uniformidad expuso una limitación de diseño de la propia métrica —su fuerte dependencia de la escala de preprocesamiento de amplitud—, que se documenta y discute en detalle. Esto refuerza que la frecuencia dominante por sí sola, sin una métrica de separabilidad correctamente normalizada, puede llevar a conclusiones poco robustas, y que el reducido número de muestras utilizadas no permite establecer conclusiones estadísticas generalizables.

\end{abstract}

% ============================================================
% PALABRAS CLAVE
% ============================================================

\begin{IEEEkeywords}

Transformada Rápida de Fourier, FFT, procesamiento digital de señales, análisis espectral, señales de audio, frecuencia dominante, separabilidad estadística.

\end{IEEEkeywords}
